% Options for packages loaded elsewhere
\PassOptionsToPackage{unicode}{hyperref}
\PassOptionsToPackage{hyphens}{url}
%
\documentclass[
]{article}
\usepackage{amsmath,amssymb}
\usepackage{iftex}
\ifPDFTeX
  \usepackage[T1]{fontenc}
  \usepackage[utf8]{inputenc}
  \usepackage{textcomp} % provide euro and other symbols
\else % if luatex or xetex
  \usepackage{unicode-math} % this also loads fontspec
  \defaultfontfeatures{Scale=MatchLowercase}
  \defaultfontfeatures[\rmfamily]{Ligatures=TeX,Scale=1}
\fi
\usepackage{lmodern}
\ifPDFTeX\else
  % xetex/luatex font selection
\fi
% Use upquote if available, for straight quotes in verbatim environments
\IfFileExists{upquote.sty}{\usepackage{upquote}}{}
\IfFileExists{microtype.sty}{% use microtype if available
  \usepackage[]{microtype}
  \UseMicrotypeSet[protrusion]{basicmath} % disable protrusion for tt fonts
}{}
\makeatletter
\@ifundefined{KOMAClassName}{% if non-KOMA class
  \IfFileExists{parskip.sty}{%
    \usepackage{parskip}
  }{% else
    \setlength{\parindent}{0pt}
    \setlength{\parskip}{6pt plus 2pt minus 1pt}}
}{% if KOMA class
  \KOMAoptions{parskip=half}}
\makeatother
\usepackage{xcolor}
\usepackage[margin=1in]{geometry}
\usepackage{color}
\usepackage{fancyvrb}
\newcommand{\VerbBar}{|}
\newcommand{\VERB}{\Verb[commandchars=\\\{\}]}
\DefineVerbatimEnvironment{Highlighting}{Verbatim}{commandchars=\\\{\}}
% Add ',fontsize=\small' for more characters per line
\usepackage{framed}
\definecolor{shadecolor}{RGB}{248,248,248}
\newenvironment{Shaded}{\begin{snugshade}}{\end{snugshade}}
\newcommand{\AlertTok}[1]{\textcolor[rgb]{0.94,0.16,0.16}{#1}}
\newcommand{\AnnotationTok}[1]{\textcolor[rgb]{0.56,0.35,0.01}{\textbf{\textit{#1}}}}
\newcommand{\AttributeTok}[1]{\textcolor[rgb]{0.13,0.29,0.53}{#1}}
\newcommand{\BaseNTok}[1]{\textcolor[rgb]{0.00,0.00,0.81}{#1}}
\newcommand{\BuiltInTok}[1]{#1}
\newcommand{\CharTok}[1]{\textcolor[rgb]{0.31,0.60,0.02}{#1}}
\newcommand{\CommentTok}[1]{\textcolor[rgb]{0.56,0.35,0.01}{\textit{#1}}}
\newcommand{\CommentVarTok}[1]{\textcolor[rgb]{0.56,0.35,0.01}{\textbf{\textit{#1}}}}
\newcommand{\ConstantTok}[1]{\textcolor[rgb]{0.56,0.35,0.01}{#1}}
\newcommand{\ControlFlowTok}[1]{\textcolor[rgb]{0.13,0.29,0.53}{\textbf{#1}}}
\newcommand{\DataTypeTok}[1]{\textcolor[rgb]{0.13,0.29,0.53}{#1}}
\newcommand{\DecValTok}[1]{\textcolor[rgb]{0.00,0.00,0.81}{#1}}
\newcommand{\DocumentationTok}[1]{\textcolor[rgb]{0.56,0.35,0.01}{\textbf{\textit{#1}}}}
\newcommand{\ErrorTok}[1]{\textcolor[rgb]{0.64,0.00,0.00}{\textbf{#1}}}
\newcommand{\ExtensionTok}[1]{#1}
\newcommand{\FloatTok}[1]{\textcolor[rgb]{0.00,0.00,0.81}{#1}}
\newcommand{\FunctionTok}[1]{\textcolor[rgb]{0.13,0.29,0.53}{\textbf{#1}}}
\newcommand{\ImportTok}[1]{#1}
\newcommand{\InformationTok}[1]{\textcolor[rgb]{0.56,0.35,0.01}{\textbf{\textit{#1}}}}
\newcommand{\KeywordTok}[1]{\textcolor[rgb]{0.13,0.29,0.53}{\textbf{#1}}}
\newcommand{\NormalTok}[1]{#1}
\newcommand{\OperatorTok}[1]{\textcolor[rgb]{0.81,0.36,0.00}{\textbf{#1}}}
\newcommand{\OtherTok}[1]{\textcolor[rgb]{0.56,0.35,0.01}{#1}}
\newcommand{\PreprocessorTok}[1]{\textcolor[rgb]{0.56,0.35,0.01}{\textit{#1}}}
\newcommand{\RegionMarkerTok}[1]{#1}
\newcommand{\SpecialCharTok}[1]{\textcolor[rgb]{0.81,0.36,0.00}{\textbf{#1}}}
\newcommand{\SpecialStringTok}[1]{\textcolor[rgb]{0.31,0.60,0.02}{#1}}
\newcommand{\StringTok}[1]{\textcolor[rgb]{0.31,0.60,0.02}{#1}}
\newcommand{\VariableTok}[1]{\textcolor[rgb]{0.00,0.00,0.00}{#1}}
\newcommand{\VerbatimStringTok}[1]{\textcolor[rgb]{0.31,0.60,0.02}{#1}}
\newcommand{\WarningTok}[1]{\textcolor[rgb]{0.56,0.35,0.01}{\textbf{\textit{#1}}}}
\usepackage{graphicx}
\makeatletter
\newsavebox\pandoc@box
\newcommand*\pandocbounded[1]{% scales image to fit in text height/width
  \sbox\pandoc@box{#1}%
  \Gscale@div\@tempa{\textheight}{\dimexpr\ht\pandoc@box+\dp\pandoc@box\relax}%
  \Gscale@div\@tempb{\linewidth}{\wd\pandoc@box}%
  \ifdim\@tempb\p@<\@tempa\p@\let\@tempa\@tempb\fi% select the smaller of both
  \ifdim\@tempa\p@<\p@\scalebox{\@tempa}{\usebox\pandoc@box}%
  \else\usebox{\pandoc@box}%
  \fi%
}
% Set default figure placement to htbp
\def\fps@figure{htbp}
\makeatother
\setlength{\emergencystretch}{3em} % prevent overfull lines
\providecommand{\tightlist}{%
  \setlength{\itemsep}{0pt}\setlength{\parskip}{0pt}}
\setcounter{secnumdepth}{-\maxdimen} % remove section numbering
\usepackage{graphicx}
\usepackage{float}
\usepackage{amsmath}
\usepackage{bookmark}
\IfFileExists{xurl.sty}{\usepackage{xurl}}{} % add URL line breaks if available
\urlstyle{same}
\hypersetup{
  hidelinks,
  pdfcreator={LaTeX via pandoc}}

\author{}
\date{\vspace{-2.5em}}

\begin{document}

\begin{Shaded}
\begin{Highlighting}[]
\FunctionTok{options}\NormalTok{(}\AttributeTok{OutDec =} \StringTok{"."}\NormalTok{)  }\CommentTok{\# Asegurar punto decimal en este chunk}

\CommentTok{\# Establecer semilla aleatoria}
\FunctionTok{set.seed}\NormalTok{(}\FunctionTok{sample}\NormalTok{(}\DecValTok{1}\SpecialCharTok{:}\DecValTok{10000}\NormalTok{, }\DecValTok{1}\NormalTok{))}

\CommentTok{\# Aleatorización del precio original}
\NormalTok{precios\_base }\OtherTok{\textless{}{-}} \FunctionTok{c}\NormalTok{(}\DecValTok{50000}\NormalTok{, }\DecValTok{60000}\NormalTok{, }\DecValTok{75000}\NormalTok{, }\DecValTok{80000}\NormalTok{, }\DecValTok{90000}\NormalTok{, }\DecValTok{100000}\NormalTok{, }\DecValTok{120000}\NormalTok{, }\DecValTok{125000}\NormalTok{,}
                 \DecValTok{150000}\NormalTok{, }\DecValTok{180000}\NormalTok{, }\DecValTok{200000}\NormalTok{, }\DecValTok{250000}\NormalTok{, }\DecValTok{300000}\NormalTok{, }\DecValTok{350000}\NormalTok{, }\DecValTok{400000}\NormalTok{, }\DecValTok{450000}\NormalTok{, }\DecValTok{500000}\NormalTok{)}
\NormalTok{precio\_original }\OtherTok{\textless{}{-}} \FunctionTok{sample}\NormalTok{(precios\_base, }\DecValTok{1}\NormalTok{)}

\CommentTok{\# Aleatorización del porcentaje de descuento}
\NormalTok{porcentajes }\OtherTok{\textless{}{-}} \FunctionTok{c}\NormalTok{(}\DecValTok{10}\NormalTok{, }\DecValTok{15}\NormalTok{, }\DecValTok{20}\NormalTok{, }\DecValTok{25}\NormalTok{, }\DecValTok{30}\NormalTok{, }\DecValTok{33}\NormalTok{, }\DecValTok{40}\NormalTok{, }\DecValTok{50}\NormalTok{)}
\NormalTok{porcentaje\_descuento }\OtherTok{\textless{}{-}} \FunctionTok{sample}\NormalTok{(porcentajes, }\DecValTok{1}\NormalTok{)}

\CommentTok{\# Aleatorización del contexto del problema}
\NormalTok{contextos\_articulos }\OtherTok{\textless{}{-}} \FunctionTok{list}\NormalTok{(}
  \FunctionTok{c}\NormalTok{(}\StringTok{"un televisor"}\NormalTok{, }\StringTok{"televisor"}\NormalTok{, }\StringTok{"electrónico"}\NormalTok{),}
  \FunctionTok{c}\NormalTok{(}\StringTok{"un celular"}\NormalTok{, }\StringTok{"celular"}\NormalTok{, }\StringTok{"electrónico"}\NormalTok{),}
  \FunctionTok{c}\NormalTok{(}\StringTok{"un computador"}\NormalTok{, }\StringTok{"computador"}\NormalTok{, }\StringTok{"electrónico"}\NormalTok{),}
  \FunctionTok{c}\NormalTok{(}\StringTok{"una lavadora"}\NormalTok{, }\StringTok{"lavadora"}\NormalTok{, }\StringTok{"electrodoméstico"}\NormalTok{),}
  \FunctionTok{c}\NormalTok{(}\StringTok{"un refrigerador"}\NormalTok{, }\StringTok{"refrigerador"}\NormalTok{, }\StringTok{"electrodoméstico"}\NormalTok{),}
  \FunctionTok{c}\NormalTok{(}\StringTok{"una bicicleta"}\NormalTok{, }\StringTok{"bicicleta"}\NormalTok{, }\StringTok{"artículo deportivo"}\NormalTok{),}
  \FunctionTok{c}\NormalTok{(}\StringTok{"un sofá"}\NormalTok{, }\StringTok{"sofá"}\NormalTok{, }\StringTok{"mueble"}\NormalTok{),}
  \FunctionTok{c}\NormalTok{(}\StringTok{"una mesa"}\NormalTok{, }\StringTok{"mesa"}\NormalTok{, }\StringTok{"mueble"}\NormalTok{),}
  \FunctionTok{c}\NormalTok{(}\StringTok{"un vestido"}\NormalTok{, }\StringTok{"vestido"}\NormalTok{, }\StringTok{"prenda"}\NormalTok{),}
  \FunctionTok{c}\NormalTok{(}\StringTok{"un traje"}\NormalTok{, }\StringTok{"traje"}\NormalTok{, }\StringTok{"prenda"}\NormalTok{),}
  \FunctionTok{c}\NormalTok{(}\StringTok{"unos zapatos"}\NormalTok{, }\StringTok{"zapatos"}\NormalTok{, }\StringTok{"calzado"}\NormalTok{),}
  \FunctionTok{c}\NormalTok{(}\StringTok{"un reloj"}\NormalTok{, }\StringTok{"reloj"}\NormalTok{, }\StringTok{"accesorio"}\NormalTok{),}
  \FunctionTok{c}\NormalTok{(}\StringTok{"una joya"}\NormalTok{, }\StringTok{"joya"}\NormalTok{, }\StringTok{"accesorio"}\NormalTok{),}
  \FunctionTok{c}\NormalTok{(}\StringTok{"un libro"}\NormalTok{, }\StringTok{"libro"}\NormalTok{, }\StringTok{"artículo educativo"}\NormalTok{),}
  \FunctionTok{c}\NormalTok{(}\StringTok{"un boleto de avión"}\NormalTok{, }\StringTok{"boleto"}\NormalTok{, }\StringTok{"servicio"}\NormalTok{)}
\NormalTok{)}
\NormalTok{contexto\_seleccionado }\OtherTok{\textless{}{-}} \FunctionTok{sample}\NormalTok{(contextos\_articulos, }\DecValTok{1}\NormalTok{)[[}\DecValTok{1}\NormalTok{]]}
\NormalTok{articulo }\OtherTok{\textless{}{-}}\NormalTok{ contexto\_seleccionado[}\DecValTok{1}\NormalTok{]}
\NormalTok{tipo\_articulo }\OtherTok{\textless{}{-}}\NormalTok{ contexto\_seleccionado[}\DecValTok{2}\NormalTok{]}
\NormalTok{categoria }\OtherTok{\textless{}{-}}\NormalTok{ contexto\_seleccionado[}\DecValTok{3}\NormalTok{]}

\CommentTok{\# Aleatorización de términos para el enunciado}
\NormalTok{terminos\_ahorro }\OtherTok{\textless{}{-}} \FunctionTok{c}\NormalTok{(}\StringTok{"ahorró"}\NormalTok{, }\StringTok{"economizó"}\NormalTok{, }\StringTok{"se descontó"}\NormalTok{, }\StringTok{"dejó de pagar"}\NormalTok{, }\StringTok{"redujo del pago"}\NormalTok{)}
\NormalTok{termino\_ahorro }\OtherTok{\textless{}{-}} \FunctionTok{sample}\NormalTok{(terminos\_ahorro, }\DecValTok{1}\NormalTok{)}

\NormalTok{terminos\_descuento }\OtherTok{\textless{}{-}} \FunctionTok{c}\NormalTok{(}\StringTok{"descuento"}\NormalTok{, }\StringTok{"rebaja"}\NormalTok{, }\StringTok{"reducción"}\NormalTok{, }\StringTok{"oferta"}\NormalTok{, }\StringTok{"promoción"}\NormalTok{)}
\NormalTok{termino\_descuento }\OtherTok{\textless{}{-}} \FunctionTok{sample}\NormalTok{(terminos\_descuento, }\DecValTok{1}\NormalTok{)}

\NormalTok{terminos\_compra }\OtherTok{\textless{}{-}} \FunctionTok{c}\NormalTok{(}\StringTok{"comprar"}\NormalTok{, }\StringTok{"adquirir"}\NormalTok{, }\StringTok{"obtener"}\NormalTok{, }\StringTok{"conseguir"}\NormalTok{)}
\NormalTok{termino\_compra }\OtherTok{\textless{}{-}} \FunctionTok{sample}\NormalTok{(terminos\_compra, }\DecValTok{1}\NormalTok{)}

\NormalTok{terminos\_costo }\OtherTok{\textless{}{-}} \FunctionTok{c}\NormalTok{(}\StringTok{"costaba"}\NormalTok{, }\StringTok{"tenía un precio de"}\NormalTok{, }\StringTok{"valía"}\NormalTok{, }\StringTok{"estaba marcado en"}\NormalTok{)}
\NormalTok{termino\_costo }\OtherTok{\textless{}{-}} \FunctionTok{sample}\NormalTok{(terminos\_costo, }\DecValTok{1}\NormalTok{)}

\NormalTok{terminos\_procedimiento }\OtherTok{\textless{}{-}} \FunctionTok{c}\NormalTok{(}\StringTok{"procedimientos"}\NormalTok{, }\StringTok{"métodos"}\NormalTok{, }\StringTok{"cálculos"}\NormalTok{, }\StringTok{"operaciones"}\NormalTok{, }\StringTok{"fórmulas"}\NormalTok{)}
\NormalTok{termino\_procedimiento }\OtherTok{\textless{}{-}} \FunctionTok{sample}\NormalTok{(terminos\_procedimiento, }\DecValTok{1}\NormalTok{)}

\CommentTok{\# Cálculos para las opciones}
\CommentTok{\# Valor del descuento (ahorro)}
\NormalTok{valor\_descuento }\OtherTok{\textless{}{-}}\NormalTok{ precio\_original }\SpecialCharTok{*}\NormalTok{ porcentaje\_descuento }\SpecialCharTok{/} \DecValTok{100}
\CommentTok{\# Precio final después del descuento}
\NormalTok{precio\_final }\OtherTok{\textless{}{-}}\NormalTok{ precio\_original }\SpecialCharTok{{-}}\NormalTok{ valor\_descuento}
\CommentTok{\# Factor para el precio final (1 {-} porcentaje\_descuento/100)}
\NormalTok{factor\_precio\_final }\OtherTok{\textless{}{-}} \DecValTok{1} \SpecialCharTok{{-}}\NormalTok{ porcentaje\_descuento}\SpecialCharTok{/}\DecValTok{100}

\CommentTok{\# Generar las opciones de respuesta}
\CommentTok{\# Opción A: Fracción equivalente al porcentaje × precio}
\NormalTok{fraccion }\OtherTok{\textless{}{-}} \ControlFlowTok{switch}\NormalTok{(}\FunctionTok{as.character}\NormalTok{(porcentaje\_descuento),}
                  \StringTok{"10"} \OtherTok{=} \StringTok{"}\SpecialCharTok{\textbackslash{}\textbackslash{}}\StringTok{frac\{1\}\{10\}"}\NormalTok{,}
                  \StringTok{"15"} \OtherTok{=} \StringTok{"}\SpecialCharTok{\textbackslash{}\textbackslash{}}\StringTok{frac\{3\}\{20\}"}\NormalTok{,}
                  \StringTok{"20"} \OtherTok{=} \StringTok{"}\SpecialCharTok{\textbackslash{}\textbackslash{}}\StringTok{frac\{1\}\{5\}"}\NormalTok{,}
                  \StringTok{"25"} \OtherTok{=} \StringTok{"}\SpecialCharTok{\textbackslash{}\textbackslash{}}\StringTok{frac\{1\}\{4\}"}\NormalTok{,}
                  \StringTok{"30"} \OtherTok{=} \StringTok{"}\SpecialCharTok{\textbackslash{}\textbackslash{}}\StringTok{frac\{3\}\{10\}"}\NormalTok{,}
                  \StringTok{"33"} \OtherTok{=} \StringTok{"}\SpecialCharTok{\textbackslash{}\textbackslash{}}\StringTok{frac\{1\}\{3\}"}\NormalTok{,}
                  \StringTok{"40"} \OtherTok{=} \StringTok{"}\SpecialCharTok{\textbackslash{}\textbackslash{}}\StringTok{frac\{2\}\{5\}"}\NormalTok{,}
                  \StringTok{"50"} \OtherTok{=} \StringTok{"}\SpecialCharTok{\textbackslash{}\textbackslash{}}\StringTok{frac\{1\}\{2\}"}\NormalTok{)}

\NormalTok{opcion\_a }\OtherTok{\textless{}{-}} \FunctionTok{paste0}\NormalTok{(fraccion, }\StringTok{" }\SpecialCharTok{\textbackslash{}\textbackslash{}}\StringTok{times "}\NormalTok{, }\FunctionTok{format}\NormalTok{(precio\_original, }\AttributeTok{big.mark=}\StringTok{"."}\NormalTok{, }\AttributeTok{decimal.mark=}\StringTok{","}\NormalTok{))}

\CommentTok{\# Opción B: Decimal equivalente al porcentaje × precio}
\NormalTok{decimal }\OtherTok{\textless{}{-}}\NormalTok{ porcentaje\_descuento }\SpecialCharTok{/} \DecValTok{100}
\NormalTok{opcion\_b }\OtherTok{\textless{}{-}} \FunctionTok{paste0}\NormalTok{(}\FunctionTok{format}\NormalTok{(decimal, }\AttributeTok{nsmall=}\DecValTok{2}\NormalTok{, }\AttributeTok{decimal.mark=}\StringTok{","}\NormalTok{), }\StringTok{" }\SpecialCharTok{\textbackslash{}\textbackslash{}}\StringTok{times "}\NormalTok{, }\FunctionTok{format}\NormalTok{(precio\_original, }\AttributeTok{big.mark=}\StringTok{"."}\NormalTok{, }\AttributeTok{decimal.mark=}\StringTok{","}\NormalTok{))}

\CommentTok{\# Opción C: Fórmula estándar de porcentaje}
\NormalTok{opcion\_c }\OtherTok{\textless{}{-}} \FunctionTok{paste0}\NormalTok{(}\StringTok{"}\SpecialCharTok{\textbackslash{}\textbackslash{}}\StringTok{frac\{"}\NormalTok{, }\FunctionTok{format}\NormalTok{(precio\_original, }\AttributeTok{big.mark=}\StringTok{"."}\NormalTok{, }\AttributeTok{decimal.mark=}\StringTok{","}\NormalTok{), }\StringTok{" }\SpecialCharTok{\textbackslash{}\textbackslash{}}\StringTok{times "}\NormalTok{, porcentaje\_descuento, }\StringTok{"\}\{100\}"}\NormalTok{)}

\CommentTok{\# Opción D: Factor para precio final × precio (NO calcula el ahorro)}
\NormalTok{opcion\_d }\OtherTok{\textless{}{-}} \FunctionTok{paste0}\NormalTok{(}\FunctionTok{format}\NormalTok{(factor\_precio\_final, }\AttributeTok{nsmall=}\DecValTok{2}\NormalTok{, }\AttributeTok{decimal.mark=}\StringTok{","}\NormalTok{), }\StringTok{" }\SpecialCharTok{\textbackslash{}\textbackslash{}}\StringTok{times "}\NormalTok{, }\FunctionTok{format}\NormalTok{(precio\_original, }\AttributeTok{big.mark=}\StringTok{"."}\NormalTok{, }\AttributeTok{decimal.mark=}\StringTok{","}\NormalTok{))}

\CommentTok{\# Crear vector con todas las opciones}
\NormalTok{opciones }\OtherTok{\textless{}{-}} \FunctionTok{c}\NormalTok{(opcion\_a, opcion\_b, opcion\_c, opcion\_d)}

\CommentTok{\# La respuesta correcta es la opción D (que NO calcula el ahorro)}
\NormalTok{indice\_correcto }\OtherTok{\textless{}{-}} \DecValTok{4}

\CommentTok{\# Crear el vector de solución para r{-}exams (1 para la respuesta correcta, 0 para las incorrectas)}
\NormalTok{solucion }\OtherTok{\textless{}{-}} \FunctionTok{rep}\NormalTok{(}\DecValTok{0}\NormalTok{, }\DecValTok{4}\NormalTok{)}
\NormalTok{solucion[indice\_correcto] }\OtherTok{\textless{}{-}} \DecValTok{1}

\CommentTok{\# Formatear valores para Python}
\NormalTok{precio\_original\_formateado }\OtherTok{\textless{}{-}} \FunctionTok{gsub}\NormalTok{(}\StringTok{"}\SpecialCharTok{\textbackslash{}\textbackslash{}}\StringTok{."}\NormalTok{, }\StringTok{","}\NormalTok{, }\FunctionTok{format}\NormalTok{(precio\_original, }\AttributeTok{big.mark=}\StringTok{"."}\NormalTok{))}
\NormalTok{valor\_descuento\_formateado }\OtherTok{\textless{}{-}} \FunctionTok{gsub}\NormalTok{(}\StringTok{"}\SpecialCharTok{\textbackslash{}\textbackslash{}}\StringTok{."}\NormalTok{, }\StringTok{","}\NormalTok{, }\FunctionTok{format}\NormalTok{(}\FunctionTok{round}\NormalTok{(valor\_descuento), }\AttributeTok{big.mark=}\StringTok{"."}\NormalTok{))}
\end{Highlighting}
\end{Shaded}

\begin{Shaded}
\begin{Highlighting}[]
\FunctionTok{options}\NormalTok{(}\AttributeTok{OutDec =} \StringTok{"."}\NormalTok{)  }\CommentTok{\# Asegurar punto decimal en este chunk}

\CommentTok{\# Crear un archivo Python temporal con el código}
\NormalTok{archivo\_python }\OtherTok{\textless{}{-}} \FunctionTok{tempfile}\NormalTok{(}\AttributeTok{fileext =} \StringTok{".py"}\NormalTok{)}

\CommentTok{\# Preparar el código Python sin usar sprintf para evitar problemas con los marcadores \%}
\NormalTok{codigo\_python }\OtherTok{\textless{}{-}} \FunctionTok{paste0}\NormalTok{(}\StringTok{\textquotesingle{}}
\StringTok{import matplotlib}
\StringTok{matplotlib.use("Agg")  \# Usar backend no interactivo}
\StringTok{import matplotlib.pyplot as plt}
\StringTok{import numpy as np}
\StringTok{from matplotlib.patches import Rectangle}
\StringTok{import matplotlib.patheffects as path\_effects}

\StringTok{\# Configuración de colores}
\StringTok{color\_precio\_original = "\#3498db"  \# Azul}
\StringTok{color\_descuento = "\#e74c3c"        \# Rojo}
\StringTok{color\_precio\_final = "\#2ecc71"     \# Verde}

\StringTok{\# Crear figura}
\StringTok{fig, ax = plt.subplots(figsize=(8, 4))}

\StringTok{\# Datos para la visualización}
\StringTok{precio\_original = \textquotesingle{}}\NormalTok{, precio\_original, }\StringTok{\textquotesingle{}}
\StringTok{porcentaje\_descuento = \textquotesingle{}}\NormalTok{, porcentaje\_descuento, }\StringTok{\textquotesingle{}}
\StringTok{valor\_descuento = precio\_original * porcentaje\_descuento / 100}
\StringTok{precio\_final = precio\_original {-} valor\_descuento}

\StringTok{\# Crear barras para representar los precios}
\StringTok{bar\_width = 0.4}
\StringTok{x\_pos = [0, 1]}

\StringTok{\# Barra para precio original}
\StringTok{ax.bar(x\_pos[0], precio\_original, bar\_width, color=color\_precio\_original,}
\StringTok{       edgecolor="black", linewidth=1.5, label="Precio original")}

\StringTok{\# Barra para precio final y descuento}
\StringTok{ax.bar(x\_pos[1], precio\_final, bar\_width, color=color\_precio\_final,}
\StringTok{       edgecolor="black", linewidth=1.5, label="Precio final")}
\StringTok{ax.bar(x\_pos[1], valor\_descuento, bar\_width, bottom=precio\_final,}
\StringTok{       color=color\_descuento, edgecolor="black", linewidth=1.5,}
\StringTok{       hatch="///", label="Descuento (\textquotesingle{}}\NormalTok{, porcentaje\_descuento, }\StringTok{\textquotesingle{}\%)")}

\StringTok{\# Añadir etiquetas con valores}
\StringTok{def add\_value\_label(x, y, value, color="black", fontsize=10, fontweight="bold"):}
\StringTok{    text = ax.text(x, y, f"$\{value:,.0f\}".replace(",", "."),}
\StringTok{                  ha="center", va="bottom", fontsize=fontsize, fontweight=fontweight)}
\StringTok{    text.set\_path\_effects([path\_effects.withStroke(linewidth=3, foreground="white")])}
\StringTok{    return text}

\StringTok{\# Etiqueta para precio original}
\StringTok{add\_value\_label(x\_pos[0], precio\_original, precio\_original)}

\StringTok{\# Etiqueta para precio final}
\StringTok{add\_value\_label(x\_pos[1], precio\_final/2, precio\_final)}

\StringTok{\# Etiqueta para descuento}
\StringTok{add\_value\_label(x\_pos[1], precio\_final + valor\_descuento/2, valor\_descuento)}

\StringTok{\# Añadir porcentaje de descuento}
\StringTok{ax.text(x\_pos[1] + 0.25, precio\_final + valor\_descuento/2,}
\StringTok{       f"\{porcentaje\_descuento\}\%", ha="left", va="center",}
\StringTok{       fontsize=12, fontweight="bold", color=color\_descuento)}

\StringTok{\# Configurar ejes}
\StringTok{ax.set\_xticks(x\_pos)}
\StringTok{ax.set\_xticklabels(["Precio original", "Precio con descuento"], fontsize=10, fontweight="bold")}
\StringTok{ax.set\_ylabel("Precio ($)", fontsize=10, fontweight="bold")}
\StringTok{ax.set\_title("Representación del descuento", fontsize=14, fontweight="bold", pad=15)}

\StringTok{\# Eliminar bordes superiores y derechos}
\StringTok{ax.spines["top"].set\_visible(False)}
\StringTok{ax.spines["right"].set\_visible(False)}

\StringTok{\# Añadir leyenda}
\StringTok{ax.legend(loc="upper center", bbox\_to\_anchor=(0.5, {-}0.15), ncol=3, frameon=False, fontsize=9)}

\StringTok{\# Crear una imagen separada para las fórmulas}
\StringTok{fig\_formula = plt.figure(figsize=(6, 1.2))}
\StringTok{plt.axis("off")}
\StringTok{plt.text(0.5, 0.7, "Ahorro = Precio original × Porcentaje de descuento ÷ 100",}
\StringTok{         ha="center", va="center", fontsize=10)}

\StringTok{\# Formatear los valores directamente en Python sin el símbolo $}
\StringTok{precio\_formateado = "\{:,.0f\}".format(precio\_original).replace(",", ".")}
\StringTok{descuento\_formateado = "\{:,.0f\}".format(valor\_descuento).replace(",", ".")}

\StringTok{\# Usar formato simple sin símbolos $ para evitar problemas}
\StringTok{ejemplo\_texto = "Ejemplo: " + precio\_formateado + " × " + str(porcentaje\_descuento) + "\% ÷ 100 = " + descuento\_formateado}
\StringTok{plt.text(0.5, 0.3, ejemplo\_texto,}
\StringTok{         ha="center", va="center", fontsize=10)}

\StringTok{plt.tight\_layout()}
\StringTok{plt.savefig("formula\_descuento.png", dpi=150, bbox\_inches="tight", transparent=True)}
\StringTok{plt.close(fig\_formula)}

\StringTok{\# Ajustar diseño}
\StringTok{plt.tight\_layout(rect=[0, 0.05, 1, 0.95])}

\StringTok{\# Guardar en múltiples formatos para asegurar compatibilidad}
\StringTok{plt.savefig("visualizacion\_descuento.png", dpi=150, bbox\_inches="tight", format="png")}
\StringTok{plt.savefig("visualizacion\_descuento.pdf", dpi=150, bbox\_inches="tight", format="pdf")}
\StringTok{plt.close()}
\StringTok{\textquotesingle{}}\NormalTok{)}

\CommentTok{\# Escribir el código Python en el archivo temporal}
\FunctionTok{writeLines}\NormalTok{(codigo\_python, archivo\_python)}

\CommentTok{\# Ejecutar el archivo Python}
\FunctionTok{py\_run\_file}\NormalTok{(archivo\_python)}

\CommentTok{\# Eliminar el archivo temporal}
\FunctionTok{unlink}\NormalTok{(archivo\_python)}
\end{Highlighting}
\end{Shaded}

\section{Question}\label{question}

Se quiere saber cuánto dinero redujo del pago una persona al conseguir
unos zapatos que valía \$80.000 y tenía un promoción del 33\%.

\includegraphics[width=0.8\linewidth,height=\textheight,keepaspectratio]{visualizacion_descuento.png}

\includegraphics[width=0.8\linewidth,height=\textheight,keepaspectratio]{formula_descuento.png}

¿Cuál de los siguientes métodos \textbf{NO} permite calcular este valor?

\subsection{Answerlist}\label{answerlist}

\begin{itemize}
\tightlist
\item
  \(\frac{1}{3} \times 80.000\)
\item
  \(0,33 \times 80.000\)
\item
  \(\frac{80.000 \times 33}{100}\)
\item
  \(0,67 \times 80.000\)
\end{itemize}

\section{Solution}\label{solution}

Para resolver este problema, debemos analizar cada uno de los
procedimientos propuestos y determinar cuál \textbf{NO} permite calcular
el ahorro obtenido al aplicar un descuento del 33\% a un artículo que
cuesta \$80.000.

\subsubsection{Análisis de cada
opción:}\label{anuxe1lisis-de-cada-opciuxf3n}

\paragraph{\texorpdfstring{1.
\(\frac{1}{3} \times 80.000\).}{1. \textbackslash frac\{1\}\{3\} \textbackslash times 80.000.}}\label{frac13-times-80.000.}

Esta opción representa la fracción equivalente al porcentaje de
descuento multiplicada por el precio original:

\begin{itemize}
\tightlist
\item
  \(\frac{1}{3}\) es equivalente a 33\%
\item
  Al multiplicar esta fracción por el precio original, obtenemos el
  valor del descuento: \(\frac{1}{3} \times 80.000 = 26.400\)
\item
  \textbf{Este procedimiento SÍ permite calcular el ahorro.}
\end{itemize}

\paragraph{\texorpdfstring{2.
\(0,33 \times 80.000\)}{2. 0,33 \textbackslash times 80.000}}\label{times-80.000}

Esta opción representa el porcentaje de descuento expresado como decimal
multiplicado por el precio original:

\begin{itemize}
\item
  \(0,33\) es equivalente a 33\%
\item
  Al multiplicar este decimal por el precio original, obtenemos el valor
  del descuento:

  \(0,33 \times 80.000 = 26.400\)
\item
  \textbf{Este procedimiento SÍ permite calcular el ahorro.}
\end{itemize}

\paragraph{\texorpdfstring{3.
\(\frac{80.000 \times 33}{100}\)}{3. \textbackslash frac\{80.000 \textbackslash times 33\}\{100\}}}\label{frac80.000-times-33100}

Esta opción representa la fórmula estándar para calcular un porcentaje
del precio original:

\begin{itemize}
\tightlist
\item
  \(\frac{80.000 \times 33}{100} = 26.400\)
\item
  \textbf{Este procedimiento SÍ permite calcular el ahorro.}
\end{itemize}

\paragraph{\texorpdfstring{4.
\(0,67 \times 80.000\)}{4. 0,67 \textbackslash times 80.000}}\label{times-80.000-1}

Esta opción representa el factor para calcular el precio final (después
del descuento) multiplicado por el precio original:

\begin{itemize}
\tightlist
\item
  \(0,67\) es equivalente a \((1 - 33/100)\)
\item
  Al multiplicar este factor por el precio original, obtenemos el precio
  final después del descuento: \(0,67 \times 80.000 = 53.600\)
\item
  Este cálculo nos da el precio final después del descuento, NO el
  ahorro.
\item
  \textbf{Este procedimiento NO permite calcular el ahorro.}
\end{itemize}

\subsubsection{Verificación
matemática:}\label{verificaciuxf3n-matemuxe1tica}

\begin{itemize}
\tightlist
\item
  Precio original: \$80.000
\item
  Porcentaje de descuento: 33\%
\item
  Valor del descuento (ahorro): \$26.400 (calculado como 33\% de
  \$80.000)
\item
  Precio final: \$53.600 (calculado como \$80.000 - \$26.400)
\end{itemize}

\subsubsection{Conclusión:}\label{conclusiuxf3n}

La opción (\(0,67 \times 80.000\)) es la única que \textbf{NO} permite
calcular el ahorro, ya que calcula el precio final después del descuento
en lugar del monto ahorrado.

\subsection{Answerlist}\label{answerlist-1}

\begin{itemize}
\tightlist
\item
  Falso
\item
  Falso
\item
  Falso
\item
  Verdadero
\end{itemize}

\section{Meta-information}\label{meta-information}

exname: descuentos\_porcentajes extype: schoice exsolution: 0001
exshuffle: TRUE exsection:
Aritmética\textbar Porcentajes\textbar Descuentos

\end{document}
